% Generated draft: requires mathematical and source review.
\documentclass[11pt]{article}
\usepackage[margin=25mm]{geometry}
\usepackage{fontspec}
\setmainfont{Latin Modern Roman}
\usepackage{amsmath,amssymb,mathtools,unicode-math}
\setmathfont{Latin Modern Math}
\usepackage{enumitem,needspace,xurl,hyperref}
\setlength{\parindent}{0pt}
\setlength{\parskip}{5pt}
\setlength{\emergencystretch}{3em}
\newcommand{\N}{\mathbb{N}}
\newcommand{\Z}{\mathbb{Z}}
\newcommand{\Q}{\mathbb{Q}}
\newcommand{\R}{\mathbb{R}}
\newcommand{\C}{\mathbb{C}}
\newcommand{\F}{\mathbb{F}}
\newcommand{\A}{\mathbb{A}}
\newcommand{\PP}{\mathbb{P}}
\newcommand{\E}{\mathbb{E}}
\newcommand{\eps}{\varepsilon}
\newcommand{\dd}{\,\mathrm{d}}
\begin{document}
\section*{UnsolvedMath ID: 700672}
\textbf{Source problem code:} catalogue entry
\subsection{Definitions and mathematical statement}
All groups considered are finite. A \emph{formation} is a nonempty class of finite groups, closed under isomorphisms, homomorphic images, and finite subdirect products. Here a subgroup \(H\leq G_1\times\cdots\times G_n\) is a subdirect product if each coordinate projection maps \(H\) onto \(G_i\).

Define
\[
\mathcal A=\{G: G\text{ is a finite group and every Sylow subgroup of }G\text{ is abelian}\}.
\]
The supplied source calls this the \(A\)-formation. To interpret the requested minimality, call a formation an \(A\)-formation when it is contained in \(\mathcal A\), and order formations by inclusion. Under this interpretation, a \emph{minimal non-\(A\)-formation} is a formation \(\mathcal F\) such that
\[
\mathcal F\nsubseteq\mathcal A,
\qquad
\text{for every formation }\mathcal E\subsetneq\mathcal F,
\quad \mathcal E\subseteq\mathcal A.
\]
Equivalently, \(\mathcal F\) contains a group with a nonabelian Sylow subgroup, whereas every group in every proper subformation of \(\mathcal F\) has abelian Sylow subgroups.

\textbf{Problem (A. N. Skiba, 1986).} Classify all minimal non-\(A\)-formations.

\textit{Formulation qualification.} The source defines the formation of all finite groups with abelian Sylow subgroups, but does not explicitly define the adjective \(A\)-formation for other formations or specify the order underlying minimality. The displayed classification target uses containment in \(\mathcal A\) and inclusion of formations; these interpretive choices require confirmation against the original problem.
\subsection{Short English statement}
Under the stated inclusion interpretation, classify the formations containing a finite group with a nonabelian Sylow subgroup for which every proper subformation consists entirely of groups with abelian Sylow subgroups. The source's terminology and minimality convention require confirmation.
\subsection{Sources}
\begin{itemize}
\item Local UnsolvedMath registry, supplied complete source record: permanent ID \(700672\), exact source problem code \texttt{catalogue entry}, title \emph{Kourovka Notebook Problem 10.57}, upstream revision \texttt{d371ee917ac9647fbb06cd6fc529deec6f3dda82}.
\item \emph{The Kourovka Notebook (Unsolved Problems in Group Theory)}, original problem 10.57, 10th issue (1986); proposed by A. N. Skiba. The supplied background also identifies the 21st issue edition (2026), arXiv:1401.0300v48. These bibliographic details are supplied attribution and have not been independently verified.
\item The supplied collection description names E. I. Khukhro and V. D. Mazurov as editors. The registry labels the problem \texttt{open}; this is a supplied status label, not an independently verified assessment of the literature.
\end{itemize}
\end{document}
